\section{HTTPS, resolución local y firewall}

\subsection{Let's Encrypt}

El equipo utilizará un dominio real y validación DNS-01. Let's Encrypt comprobará un registro TXT bajo \texttt{\_acme-challenge}. Este método funciona cuando el servidor web no se expone públicamente \cite{letsencrypt-challenges}.

Después de emitir el certificado, openSUSE resolverá el dominio hacia \texttt{10.50.5.20} mediante \texttt{/etc/hosts}. El navegador debe mostrar un certificado válido cuyo nombre coincida con el dominio.

\subsection{UFW}

\begin{table}[H]
  \centering
  \begin{tabularx}{\linewidth}{l l l Y}
    \toprule
    \textbf{Servidor} & \textbf{Puerto} & \textbf{Protocolo} & \textbf{Origen} \\
    \midrule
    Web & 22222 & TCP & 10.50.5.10 \\
    Web & 80 y 443 & TCP & 10.50.5.10 \\
    DB & 22222 & TCP & 10.50.5.10 \\
    DB & 13306 & TCP & 10.50.5.20 \\
    \bottomrule
  \end{tabularx}
  \caption{Reglas entrantes planificadas.}
\end{table}

La prueba negativa debe confirmar que openSUSE no alcanza MySQL directamente. La VM web sí debe conectarse a \texttt{10.50.5.30:13306}.

\evidencia{Certificado válido de Let's Encrypt}{05-seguridad/E15-certificado-letsencrypt.png}{Debe mostrar nombre, emisor y vigencia sin exponer credenciales.}
